\documentclass[11pt]{article}
\usepackage[a4paper,margin=25mm]{geometry}
\usepackage{amsmath,amssymb,amsthm,hyperref}
\newtheorem{theorem}{Theorem}
\title{A unit-coefficient counterexample to 00000001198}
\author{Ziang Ni (\texttt{ziangni-sys})}
\date{10 October 2026}
\begin{document}
\maketitle
The conjecture says that Zonal Littlewood--Richardson coefficients are
nonnegative alternating sums of even integers. Its English and Chinese
statements do not exclude the empty partition. That coefficient alone
contradicts the claimed parity restriction.

\paragraph{The actual polynomial and normalization.}
Write $Z_\lambda$ for the normalized Zonal polynomials, that is, the
Jack parameter-two specialization. The empty partition $\varnothing$
has weight zero. Its normalized polynomial is
\[
 Z_\varnothing(x_1,\ldots,x_m)=1.
\]
This is the actual degree-zero symmetric polynomial, rather than an
auxiliary scalar model. For the standard Zonal normalization, the
NIST Digital Library of Mathematical Functions, equations
\href{https://dlmf.nist.gov/35.4.E4}{35.4.4} and
\href{https://dlmf.nist.gov/35.4.E6}{35.4.6}, explicitly include the zero
partition. The latter identity at degree zero has exactly one summand
and reads $Z_\varnothing(T)=(\operatorname{tr}T)^0=1$.
Equivalently, in monic Jack normalization the zero partition has the
single monomial $m_\varnothing=1$ and no lower terms. The usual integral
normalization multiplies this by an empty product, also $1$.
Thus the example is independent of these standard normalization choices.

Define the LR coefficients by actual multiplication and expansion in
this homogeneous basis:
\[
 Z_\lambda Z_\mu
   =\sum_{|\nu|=|\lambda|+|\mu|}c_{\lambda\mu}^{\nu}Z_\nu.
\]
When $\lambda=\mu=\varnothing$, the only weight-zero partition is
$\nu=\varnothing$. Consequently
\[
 1=Z_\varnothing^2=c_{\varnothing\varnothing}^{\varnothing}
 Z_\varnothing,
 \qquad c_{\varnothing\varnothing}^{\varnothing}=1.
\]
The coefficient is unique: extraction of the constant monomial
coefficient from the polynomial equality gives $1=c$.

\begin{theorem}
The parity assertion in 00000001198 is false.
\end{theorem}
\begin{proof}
For any integers $a_0,\ldots,a_{r-1}$,
\[
 \sum_{j=0}^{r-1}(-1)^j(2a_j)
 =2\sum_{j=0}^{r-1}(-1)^ja_j
\]
is even. The LR coefficient above is $1$, which is odd, so it is no
such alternating sum. Requiring the sum to be nonnegative does not
remove this obstruction. Allowing arbitrary signs instead of strictly
alternating signs also preserves evenness.
\end{proof}

\paragraph{Scope and Lean verification.}
This disproves the stated even-integer clause, without asserting that
Zonal LR positivity or any corrected matching interpretation is false.
The source imposes no positive-weight restriction. The Lean project
uses actual multivariate polynomials over $\mathbb Q$ in an arbitrary
variable type. It constructs $Z_\varnothing$ as the zero-exponent
monomial with coefficient $1$, proves permutation invariance, proves
that every degree-zero polynomial lies in the constant basis, and
extracts the unique LR unit coefficient from multiplication. It then
proves the alternating-sum parity obstruction by induction on a list
of integers. All printed axiom audits use only the permitted standard
Lean axioms. The project pins Lean 4.19.0 and mathlib commit
\texttt{c44e0c8ee63ca166450922a373c7409c5d26b00b}.
\end{document}
